% =============================================================================
% exercises/ejercicios_cap04a.tex
% Ejercicios — Capítulo 4, Sección A
% Temas: Estructura atómica, orbitales, enlace e hidrocarburos
% =============================================================================

\section*{Ejercicios --- Secci\'on A: Estructura at\'omica e hidrocarburos}
\addcontentsline{toc}{section}{Ejercicios: Estructura at\'omica e hidrocarburos}
\label{sec:ejercicios_cap04a}

\begin{bloque_ejercicios}[Ejercicios A --- Estructura at\'omica e hidrocarburos]

\begin{ejercicios}

% ----- A1: Niveles de energía y capacidad electrónica ----------------------
\item Con base en la f\'ormula $2n^2$ del \textbf{Cuadro~\ref{tab1.cap3}}:

\begin{enumerate}[label=\alph*)]
    \item Calcule el n\'umero m\'aximo de electrones en los niveles
          $n = 1$, 2, 3 y 4.
    \item El cloro (Z~=~17) tiene la configuraci\'on
          $1s^2 2s^2 2p^6 3s^2 3p^5$. ?`Cu\'antos electrones le faltan
          al tercer nivel para completarse? ?`Qu\'e gas noble tiene
          esa configuraci\'on completa?
    \item Explique por qu\'e el valor te\'orico de 50 electrones en
          el nivel $n = 5$ nunca se ha alcanzado en ning\'un elemento
          conocido.
\end{enumerate}

\textit{[Respuesta: (a) 2, 8, 18, 32;
(b) le faltan 1 electr\'on; el Arg\'on (Ar) completa el tercer nivel
con $3s^2 3p^6$]}

% ----- A2: Subniveles y orbitales ------------------------------------------
\item Con base en el \textbf{Cuadro~\ref{tabla3.2}}:

\begin{enumerate}[label=\alph*)]
    \item ?`Cu\'antos electrones acepta cada uno de los subniveles
          $s$, $p$, $d$ y $f$? ?`Cu\'antos orbitales contiene cada uno?
    \item Indique el orden de energ\'ia creciente de los subniveles
          dentro de un nivel principal.
    \item Si cada orbital acepta m\'aximo 2 electrones (con spines
          opuestos), ?`por qu\'e el subnivel $p$ acepta 6 y no 2?
    \item Un electr\'on en el subnivel $3d$ y otro en el $4s$.
          ?`Cu\'al tiene mayor energ\'ia? Explique con el orden
          de llenado de orbitales.
\end{enumerate}

\textit{[Respuesta: (a) $s$: 2 e$^-$/1 orbital;
$p$: 6 e$^-$/3 orbitales;
$d$: 10 e$^-$/5 orbitales;
$f$: 14 e$^-$/7 orbitales]}

% ----- A3: Configuraciones electrónicas ------------------------------------
\item Usando el \textbf{Cuadro~\ref{tabla3.3}}, complete y conteste:

{%
\centering
\begin{tabular}{lccl}
\toprule
\textbf{Elemento} & \textbf{Z} & \textbf{Electrones de valencia} &
\textbf{Configuraci\'on} \\
\midrule
Ox\'igeno   & 8  & & \\
Silicio     & 14 & & \\
Potasio     & 19 & & \\
Calcio      & 20 & & \\
Magnesio    & 12 & & \\
\bottomrule
\end{tabular}
\par
}%

\begin{enumerate}[label=\alph*)]
    \item ?`Cu\'al de los cinco elementos es un metal alcalino?
          ?`Cu\'al es un gas noble? ?`Cu\'al es un no metal?
    \item Los elementos con la misma cantidad de electrones de
          valencia tienen propiedades qu\'imicas similares.
          Identifique dos elementos del cuadro que pertenezcan
          al mismo grupo de la tabla peri\'odica.
\end{enumerate}

\textit{[Respuesta: O: $1s^2 2s^2 2p^4$ (6 val.);
Si: $1s^2 2s^2 2p^6 3s^2 3p^2$ (4 val.);
K: $[\text{Ar}]4s^1$ (1 val.);
Ca: $[\text{Ar}]4s^2$ (2 val.);
Mg: $[\text{Ne}]3s^2$ (2 val.)]}

% ----- A4: Símbolos de Lewis y regla del octeto ----------------------------
\item Los s\'imbolos de Lewis representan los electrones de valencia
de un \'atomo con puntos.

\begin{enumerate}[label=\alph*)]
    \item Dibuje los s\'imbolos de Lewis para: H, C, N, O, F y Na.
    \item ?`Cu\'antos electrones no apareados (disponibles para enlace)
          tiene cada uno?
    \item Aplique la regla del octeto para explicar por qu\'e el
          ox\'igeno forma 2 enlaces covalentes y el nitr\'ogeno forma 3.
    \item Dibuje la estructura de Lewis del metano (\ce{CH4}),
          el agua (\ce{H2O}) y el amon\'iaco (\ce{NH3}), mostrando
          todos los pares de electrones compartidos y no compartidos.
\end{enumerate}

% ----- A5: Energía de ionización y afinidad electrónica --------------------
\item Con base en el \textbf{Cuadro~\ref{ionizacion}}:

\begin{enumerate}[label=\alph*)]
    \item ?`Por qu\'e la segunda energ\'ia de ionizaci\'on del sodio
          (\SI{4565}{\kilo\joule\per\mole}) es mucho mayor que la
          primera (\SI{496}{\kilo\joule\per\mole})?
    \item Compare las energ\'ias de ionizaci\'on del Li, Na y K.
          ?`C\'omo cambia la primera energ\'ia de ionizaci\'on al
          bajar en el grupo 1A? ?`Por qu\'e?
    \item La cantidad de energ\'ia absorbida o liberada cuando se
          agrega un electr\'on a un \'atomo neutro se llama:
          \underline{\hspace{4cm}}. ?`En qu\'e tipo de elementos
          (metales o no metales) suele ser mayor esta propiedad?
    \item A partir de las definiciones de energ\'ia de ionizaci\'on y
          afinidad electr\'onica, explique c\'omo se puede formar un
          i\'on negativo (ani\'on) a partir de un \'atomo neutro.
\end{enumerate}

\textit{[Respuesta: (c) Afinidad electr\'onica;
es mayor en no metales, especialmente hal\'ogenos]}

% ----- A6: Electronegatividad y tipo de enlace -----------------------------
\item La electronegatividad determina el tipo de enlace qu\'imico.

{%
\centering
\begin{tabular}{lcccl}
\toprule
\textbf{Par de \'atomos} & $\Delta$electroneg. & \textbf{Tipo de enlace} &
\textbf{Polar/No polar} \\
\midrule
H -- H       & 0    & & \\
H -- Cl      & 0.9  & & \\
Na -- Cl     & 2.1  & & \\
C -- H       & 0.4  & & \\
C -- O       & 1.0  & & \\
\bottomrule
\end{tabular}
\par
}%

\begin{enumerate}[label=\alph*)]
    \item Complete la tabla seg\'un los criterios del cap\'itulo
          (diferencia = 0: covalente no polar; peque\~na: covalente
          polar; grande: i\'onico).
    \item Explique por qu\'e el \ce{CO2} es una mol\'ecula no polar
          a pesar de tener enlaces \ce{C=O} polares.
    \item ?`C\'omo deben ser las electronegatividades de dos \'atomos
          para que formen un enlace covalente no polar?
\end{enumerate}

% ----- A7: Hibridación del carbono -----------------------------------------
\item El \'atomo de carbono puede estar hibridado en $sp^3$, $sp^2$
o $sp$.

\begin{enumerate}[label=\alph*)]
    \item Indique el tipo de hibridaci\'on del carbono en cada
          compuesto y el \'angulo de enlace esperado:

{%
\centering
\begin{tabular}{lccc}
\toprule
\textbf{Compuesto} & \textbf{Hibridaci\'on} & \textbf{\'Angulo de enlace} &
\textbf{Tipo de enlace C--C} \\
\midrule
\ce{CH4} (metano)       & & & \\
\ce{CH2=CH2} (etileno)  & & & \\
\ce{CH#CH} (acetileno)  & & & \\
\ce{C6H6} (benceno)     & & & \\
\bottomrule
\end{tabular}
\par
}%

    \item Explique en t\'erminos de la hibridaci\'on por qu\'e el
          metano es tetra\'edrico con \'angulos de $109.5^\circ$.
    \item ?`C\'omo se denomina a la mezcla de orbitales $s$ y $p$?
          ?`Qu\'e ventaja ofrece el orbital h\'ibrido resultante
          respecto a los orbitales puros?
\end{enumerate}

\textit{[Respuesta: \ce{CH4}: $sp^3$, $109.5^\circ$, enlace sencillo;
\ce{CH2=CH2}: $sp^2$, $120^\circ$, doble enlace;
\ce{CH\#CH}: $sp$, $180^\circ$, triple enlace;
\ce{C6H6}: $sp^2$, $120^\circ$, deslocalizado]}

% ----- A8: Fórmulas y nomenclatura de alcanos ------------------------------
\item Usando la f\'ormula general $\ce{C_nH_{2n+2}}$ y el
\textbf{Cuadro~\ref{alcanos}}:

\begin{enumerate}[label=\alph*)]
    \item Determine la f\'ormula molecular y la f\'ormula condensada
          para el pentano ($n = 5$) y el hexadecano ($n = 16$).
    \item Escriba la f\'ormula estructural condensada del:
          \begin{enumerate}[label=\roman*)]
              \item 3-etil-5-metiloctano
              \item 2,3-dimetilpentano
              \item 4-etil-2-metilheptano
          \end{enumerate}
    \item Nombre los siguientes compuestos seg\'un las reglas UIQPA:
          \begin{enumerate}[label=\roman*)]
              \item \ce{CH3-CH(CH3)-CH2-CH3}
              \item \ce{CH3-CH2-CH(C2H5)-CH(CH3)-CH3}
          \end{enumerate}
    \item ?`En qu\'e fase f\'isica (gas, l\'iquido o s\'olido) se
          encuentran los alcanos con m\'as de 14 \'atomos de carbono?
          Consulte el \textbf{Cuadro~\ref{alcanos}}.
\end{enumerate}

\textit{[Respuesta: (a) pentano: \ce{C5H12}, \ce{CH3(CH2)3CH3};
hexadecano: \ce{C16H34};
(c-i) 2-metilbutano;
(c-ii) 3-etil-2-metilpentano]}

% ----- A9: Alquenos, alquinos e isomería -----------------------------------
\item Con base en las f\'ormulas generales y el
\textbf{Cuadro~\ref{insat:1}}:

\begin{enumerate}[label=\alph*)]
    \item ?`Cu\'antos \'atomos de hidr\'ogeno menos que el alcano
          correspondiente tienen los alquenos? ?`Y los alquinos?
    \item Escriba la f\'ormula molecular y nombre UIQPA de:
          \begin{enumerate}[label=\roman*)]
              \item el alqueno con 4 carbonos y el doble enlace
                    en posici\'on 1
              \item el alquino con 4 carbonos y el triple enlace
                    en posici\'on 2
          \end{enumerate}
    \item Escriba la f\'ormula estructural condensada del:
          \begin{enumerate}[label=\roman*)]
              \item 3-etil-4-metil-5-octeno
              \item 3-etil-7-metil-5-nonino
          \end{enumerate}
    \item El compuesto \ce{CH3-C#C-CH(CH2CH3)-CH2-CH2-CH2-CH3}
          tiene el nombre: \underline{\hspace{5cm}}.
          \textit{[Respuesta: 4-etil-2-octino]}
\end{enumerate}

% ----- A10: Hidrocarburos aromáticos y nomenclatura ------------------------
\item Con base en la estructura del benceno y la nomenclatura
de los arom\'aticos:

\begin{enumerate}[label=\alph*)]
    \item ?`Por qu\'e el benceno no participa f\'acilmente en
          reacciones de adici\'on como los alquenos t\'ipicos?
          Explique en t\'erminos de su estructura de resonancia.
    \item Identifique y nombre los siguientes compuestos:
          (i) \ce{C6H5-CH3},
          (ii) \ce{C6H5-OH},
          (iii) \ce{C6H5-NH2},
          (iv) \ce{C6H5-CH=CH2}.
    \item Dibuje y nombre los tres is\'omeros de posici\'on del
          dimetilbenceno (xileno), indicando si son orto, meta
          o para.
    \item Los HAP (hidrocarburos polic\'iclicos arom\'aticos) son
          contaminantes de preocupaci\'on ambiental. Mencione dos
          propiedades f\'isicoquimicas que expliquen por qu\'e se
          acumulan en sedimentos y suelos en lugar de en agua o aire.
\end{enumerate}

\textit{[Respuesta: (b) (i) tolueno; (ii) fenol; (iii) anilina;
(iv) estireno]}

\end{ejercicios}

\end{bloque_ejercicios}
% =============================================================================
% FIN DE ejercicios_cap04a.tex
% =============================================================================
